% Generated draft: requires mathematical and source review.
\documentclass[11pt]{article}
\usepackage[margin=25mm]{geometry}
\usepackage{fontspec}
\setmainfont{Latin Modern Roman}
\usepackage{amsmath,amssymb,mathtools,unicode-math}
\setmathfont{Latin Modern Math}
\usepackage{enumitem,needspace,xurl,hyperref}
\setlength{\parindent}{0pt}
\setlength{\parskip}{5pt}
\setlength{\emergencystretch}{3em}
\newcommand{\N}{\mathbb{N}}
\newcommand{\Z}{\mathbb{Z}}
\newcommand{\Q}{\mathbb{Q}}
\newcommand{\R}{\mathbb{R}}
\newcommand{\C}{\mathbb{C}}
\newcommand{\F}{\mathbb{F}}
\newcommand{\A}{\mathbb{A}}
\newcommand{\PP}{\mathbb{P}}
\newcommand{\E}{\mathbb{E}}
\newcommand{\eps}{\varepsilon}
\newcommand{\dd}{\,\mathrm{d}}
\begin{document}
\section*{UnsolvedMath ID: 30000522}
\textbf{Source problem code:} catalogue entry
\subsection{Definitions and mathematical statement}
Let $B$ be an anisotropic bilinear form over $F$ in the category intended by the source, and let $(F_i,B_i)_{0\le i\le h}$ be its standard splitting tower. For an extension $L/F$, write $B_L=B\otimes_F L$ for scalar extension. In the usual symmetric-bilinear setting, anisotropy means $B(v,v)\ne0$ for every nonzero vector, and the Witt index $i_W(B_L)$ measures the dimension of a maximal totally isotropic subspace, with the precise Witt-decomposition convention to be fixed in the intended characteristic.

The question is whether every extension over which $B$ becomes isotropic has an index already achieved on the standard tower:
\[
 \forall L/F:\quad
 B_L\text{ isotropic}\quad\Longrightarrow\quad
 \exists k\in\{1,\ldots,h\}:\quad
 i_W(B_L)=i_W(B_{F_k}).
\]
The choice of $k$ may depend on $B$ and $L$. There is no assertion that one fixed tower stage works for all extensions. The comparison is with the original form $B$ extended to $F_k$, not the residual form $B_k$, and does not require an embedding of $L$ into $F_k$ or conversely. No finiteness, algebraicity, or separability restriction on $L/F$ is stated.

The source does not give the characteristic of $F$, symmetry and nondegeneracy assumptions on the bilinear form, its dimension, or the construction and stopping convention of the standard splitting tower. In characteristic two, bilinear and quadratic-form conventions cannot be interchanged silently; neither a quadratic-form splitting theorem nor a particular quadric-function-field construction is imported. The intended tower and index definitions must therefore be recovered from the report before this becomes a complete universal formulation.

The assessment says the extraction omits the invariant definitions and quantifiers. The present record gives the anisotropic-form/tower context and the existential tower-stage question, but does not resolve the remaining category conventions.
\subsection{Short English statement}
The question asks whether every positive Witt index obtained after scalar extension already occurs at some stage of the form's standard splitting tower. The stage may depend on the extension, and the index is that of the original form over the tower field. Characteristic and bilinear-form conventions are missing.
\subsection{Sources}
\begin{itemize}
\item Local UnsolvedMath registry, permanent ID 30000522, code catalogue entry, ``Witt Indices along Bilinear-Form Splitting Towers.'' Proposed year 2006; proposer unspecified.
\item \emph{Quadratic Forms and Linear Algebraic Groups} (2007), Oberwolfach Reports, pp.\ 1743--1794, DOI 10.4171/owr/2006/29; \url{https://doi.org/10.4171/owr/2006/29}. The displayed report year 2007 differs from the DOI and proposed year 2006.
\item Statement status corrected\_verified; the supplied note says editorial and LaTeX cleanup preserved meaning and added no source-dependent content.
\item Literature triage checked 2026-08-21: open. It says $i_W$, scalar-extension notation, and quantifiers were not defined in the excerpt and no general equality could be verified from the available literature. This is a dated formulation assessment, not an independently checked theorem status. Created 2026-08-25; updated 2026-09-16.
\item Registry curation by Ulam AI / UnsolvedMath Contributors, CC BY 4.0; the underlying report retains its supplied provenance. Unreviewed statement-only draft.
\end{itemize}
\end{document}
